% ==================================================================
\subsection{Task 07: Query Design}
\label{sec:task07}
% ==================================================================


\subsubsection{Database Objective}

Task~07 designed executable SQL queries that answer role-specific business
questions for the Campus Space Management System.The task produced \texttt{outputs/07-query-design-G05.sql}, containing all
student queries assembled into a single executable file. Each query follows the
4-field template: business question, target user(s), explanation of usefulness,
and a parameterized T-SQL statement using \texttt{DECLARE} blocks.

\subsubsection{Agent Configuration and Improvement Process}

\paragraph{Initial approach.}
The initial design treated Task~07 as a straightforward generation task: the
agent would read the approved schema and sample data, then produce five queries
independently. A basic command and skill were written to define the output
template and SQL quality rules. While the agent successfully generated
syntactically valid queries, the output lacked depth and variety. Queries
tended to cluster around simple availability checks and repeated similar
patterns across students, resulting in low semantic coverage of the domain's
business questions.

\paragraph{Improvement: role-scoped query generation.}
After reviewing the initial output, the team identified the root cause: the
agent had no guidance on which user perspective each student should adopt.
Without role constraints, every student's queries converged toward the same
generic availability-search pattern.

The command was extended with explicit per-student arguments:

\begin{itemize}[noitemsep]
    \item \texttt{--student-name} to assign authorship to each query block
    \item \texttt{--target-users} to restrict each student's queries to a
          specific role subset (e.g.\ \texttt{facility\_manager},
          \texttt{lecturer})
    \item \texttt{--num-queries} to control output volume
    \item \texttt{--mode} to support append or overwrite when multiple
          students contributed to the same file
    \item \texttt{--business-question} to guide the agent to follow each
          student's intended scenario
\end{itemize}

The skill was updated with a query distribution table mapping each role to
expected query categories: lecturers focused on availability and schedule
views; facility staff on approvals, check-in, and maintenance; facility
managers on utilization analytics and no-show analysis; department admins on
audit trails and inter-department comparisons. This prevented semantic
duplication across students and produced a query set with wider operational
coverage.

\paragraph{Interaction rule.}
To protect each student's contribution, the command required the agent to stop
and ask the user before overwriting an existing output file. If \texttt{--mode}
was not specified and the file already existed, the agent would prompt:
\textit{``Output file exists. Append new queries for [student name] to the
existing file, or overwrite the entire file?''} Only after an explicit answer
would the agent proceed. This prevented silent data loss when multiple students
ran the command sequentially.

% ==================================================================
\subsubsection{Per-Student Query Summary}
% ==================================================================

% ----------------------------------------------------------------
% TEMPLATE for each student — copy and fill in
% ----------------------------------------------------------------

% ---- Student: Nguyen Huu Phuoc ----
\paragraph{Nguyễn Hữu Phước --- Role: Facility Staff.}
\begin{itemize}[noitemsep]

    \item \textbf{Q1 --- Pending approvals for today:} Retrieves all pending booking requests within the next 24 hours so facility staff can review and approve or reject them in a timely manner.

\item \textbf{Q2 --- Real-time occupancy snapshot:} Shows which spaces are currently checked-in or approved, who booked them, and when the session ends, enabling live occupancy monitoring across campus.

\item \textbf{Q3 --- Active maintenance tickets:} Lists all open and in-progress maintenance tickets with problem descriptions, assigned staff, and days open, helping staff prioritise repairs.

\item \textbf{Q4 --- Eligible check-in bookings:} Identifies approved bookings past their start time with no check-in, enabling proactive follow-up on no-shows or late arrivals.

\item \textbf{Q5 --- Session completion report for today:} Retrieves sessions completed on a given date with actual end times, space conditions, and check-in staff, supporting end-of-day space review.
\end{itemize}
% ---- Student: Pham Huu Nam ----
\paragraph{Phạm Hữu Nam --- Role: Facility Manager.}
\begin{itemize}[noitemsep]
    \item \textbf{Q6 --- Rejected booking audit trail:} retrieves the full
    rejection history for a specific lecturer, filtered by their natural
    key (unique email), including rejection reasons, decision timestamps,
    and the staff member who processed each decision. Supports
    investigation of possible bias or procedural issues and accountability
    reviews.

    \item \textbf{Q7 --- Spaces with upcoming bookings but no recent
    maintenance:} identifies spaces with approved or checked-in bookings
    scheduled within the next 10 days that have \emph{not} undergone any
    resolved maintenance or quality inspection within the past month (and
    are not currently under maintenance), so technicians can be
    proactively dispatched for pre-use checks on long-neglected,
    high-risk rooms.

    \item \textbf{Q8 --- Department no-show rate analysis:} aggregates
    confirmed booking and no-show counts per department for the current
    semester, computes the no-show rate percentage, and ranks departments
    from worst to best to support enforcement of stricter reservation
    penalties on departments that over-reserve and under-utilize spaces.

    \item \textbf{Q9 --- Cumulative usage hours for competing requesters:}
    identifies student requesters with overlapping pending or approved
    bookings for a contested meeting room and time slot, then calculates
    each requester's cumulative actual usage hours for the current month
    from completed booking sessions, enabling fair-access allocation that
    prioritises under-served (low-usage) users.

    \item \textbf{Q10 --- Semester room-type request summary:} aggregates
    total requests, approval rates, and unique user counts by room
    category (classroom, laboratory, meeting room, auditorium) for a
    declared semester window, to support data-driven decisions on space
    expansion or consolidation.
\end{itemize}
% ---- Student: Cao Quang Hung ----
\paragraph{Cao Quảng Hưng --- Role: Student.}
\begin{itemize}[noitemsep]
    \item \textbf{Q11 --- Available spaces with particular facilities:} Finds teaching or study spaces equipped with some particular facilities (such as Projector, WhiteBoard) that are free to book within their desired time window

    \item \textbf{Q12 --- Particular events on a date:} Find particular events (such as Seminar or Student Activity) are happening on a given date with specified capacity and organizer.

    \item \textbf{Q13 --- Alternative spaces when usual room is blocked: } Learns which room the student uses most from their booking history, checks if it is unavailable then proactively recommends up to a specified number of alternatives on the same floor. 

    \item \textbf{Q14 --- Lab availability by day of week:} Based on booking history from the past year, retrieves which days of the week have the most available project and computer laboratories

    \item \textbf{Q15 --- Department admin's upcoming approved bookings:} A particular department administrators identifies all approved upcoming booking made by members of their department, seeing who booked which space, for what purpose, and when.
\end{itemize}

% ---- Student: Tran Dinh Quoc Thang ----
\paragraph{Trần Đình Quốc Thắng --- Role: Lecturer, Teaching Assisstant.}
\begin{itemize}[noitemsep]
    \item \textbf{Q16 --- Lab availability with equipment and conflict checks:}
    finds computer or project labs for a requested tutorial slot by checking
    capacity, workstation count, optional projector availability, live room
    status, confirmed booking overlaps, and unresolved maintenance conflicts.

    \item \textbf{Q17 --- Lecturer personal booking timeline:} retrieves a
    lecturer's semester booking history across all active lifecycle statuses,
    including room details, decision notes, rejection reasons, approver
    information, and actual session timestamps.

    \item \textbf{Q18 --- Completed TA lab session history:} lists completed
    computer/project lab sessions for a teaching assistant, including actual
    duration, condition notes, check-in staff, usage notes, and a summarized
    facility list.

    \item \textbf{Q19 --- Approval lead-time analysis:} aggregates approved
    and rejected lecturer bookings by purpose, space type, and decision status
    to show minimum, average, and maximum decision lead time for planning
    future submissions.

    \item \textbf{Q20 --- Upcoming lab readiness alert:} identifies approved
    upcoming TA lab bookings that may be disrupted by non-available space
    status or overlapping unresolved maintenance, supporting early backup-room
    planning.
\end{itemize}

% ==================================================================
\subsubsection{Issues and Improvements Summary}
% ==================================================================

Table~\ref{tab:task07-improvements} consolidates the significant issues
identified across all students during the review cycle, together with the
improvements made and their effect on query correctness or maintainability.

\begin{table}[http]
\centering
\footnotesize
\renewcommand{\arraystretch}{1.4}
\begin{tabularx}{\textwidth}{>{\raggedright\arraybackslash}p{0.6cm}
                             >{\raggedright\arraybackslash}p{1.2cm}
                             >{\raggedright\arraybackslash}p{3.2cm}
                             >{\raggedright\arraybackslash}X
                             >{\raggedright\arraybackslash}p{2.6cm}}
\toprule
\textbf{Q\#} & \textbf{Student} & \textbf{Issue found} & \textbf{Improvement made} & \textbf{Effect} \\
\midrule

Q6 & Nam &
Filtered on \texttt{full\_name}, which is not unique. &
Replaced with \texttt{DECLARE @requester\_id INT}; filter now uses the
surrogate key \texttt{b.requester\_id}. &
Eliminates false multi-row matches. \\


\addlinespace[0.4em]


Q9 & Nam &
\texttt{COALESCE(actual\_end, requested\_end)} overcounted
hours for \texttt{checked\_in} sessions still in progress. &
Replaced with a \texttt{CASE} expression: \texttt{completed} uses
\texttt{actual\_end\_time}; \texttt{checked\_in} in-progress uses
\texttt{GETDATE()}; month bounds derived via \texttt{DATEFROMPARTS}. &
Duration accurate per status; query is self-updating each month. \\

\addlinespace[0.4em]

% ---- Add rows for other students below ----
% Q11 & Phuoc & [issue] & [improvement] & [effect] \\
% Q12 & Hung  & [issue] & [improvement] & [effect] \\

\bottomrule
\end{tabularx}
\caption{Issues and improvements across all Task~07 student queries}
\label{tab:task07-improvements}
\end{table}

% ==================================================================
\subsubsection{Evaluation Result}
% ==================================================================

Task~07 was evaluated on query quality, SQL correctness, and alignment with the
4-field template.
Queries that returned zero rows due to sample data coverage gaps were documented
with an inline comment noting that the logic was correct and the data seed did
not cover the scenario. No query was accepted without at least a syntax
verification pass.

\newpage




